\documentclass[11pt]{article}
\usepackage{amsmath,amssymb}
\usepackage{geometry}
\geometry{margin=1in}

\title{Equivalence Completion Between Electroweak and Electromagnetic Coupling Bookkeeping}
\author{Amos Jay Maley}
\date{}

\begin{document}
\maketitle

\section*{Abstract}
We complete an equivalence map between electroweak gauge bookkeeping and low-energy electromagnetic bookkeeping. Parameters commonly treated as independent in the ultraviolet description are shown to overcount tensor load once admissible redescription and regime transport are enforced. The electromagnetic sector admits a single invariant interaction strength. The result is eliminative, invariant under admissible redescription, and introduces no new physics.

\section{Apparent Independence in Electroweak Gauge Bookkeeping}
In electroweak gauge bookkeeping, the interaction terms are written as
\begin{equation}
\mathcal{L}_{\text{gauge}} \supset g_2\, W^a_\mu J^{a\mu} + g_1\, B_\mu J_Y^\mu ,
\end{equation}
with the weak mixing angle defined by
\begin{equation}
\tan\theta_W = \frac{g_1}{g_2}.
\end{equation}
The parameters $g_2$, $g_1$, and $\theta_W$ are often treated as independent physical inputs in this skin.

\section{Admissible Electromagnetic Bookkeeping}
An admissible low-energy description of electromagnetic interactions is given by quantum electrodynamics, in which all electromagnetic processes are governed by a single coupling $e$,
\begin{equation}
\mathcal{L}_{\text{EM}} \supset e\, A_\mu J_{\text{EM}}^\mu .
\end{equation}
No independent parameters corresponding to $g_2$, $g_1$, or $\theta_W$ appear in this skin.

\section{Lemma 1: Cross-Skin Transport}
Electroweak symmetry breaking induces a field rotation
\begin{equation}
\begin{pmatrix}
A_\mu \\
Z_\mu
\end{pmatrix}
=
\begin{pmatrix}
\cos\theta_W & \sin\theta_W \\
-\sin\theta_W & \cos\theta_W
\end{pmatrix}
\begin{pmatrix}
B_\mu \\
W^3_\mu
\end{pmatrix}.
\end{equation}
Under this admissible redescription, the electromagnetic coupling is identified as
\begin{equation}
e = g_2 \sin\theta_W = g_1 \cos\theta_W .
\end{equation}

\section{Lemma 2: Invariant Witness Counting}
If a tensor load contains $n$ independent degrees of freedom, admissible redescription and regime transport must preserve $n$ independent invariant witnesses across all admissible skins describing the same regime. A skin exposing fewer invariant witnesses implies redundancy in higher-dimensional parameterizations.

\section{Failure of Independence}
Assume $(g_2, g_1)$ constitute independent tensor degrees of freedom for electromagnetic interaction strength. By Lemma~2, two independent invariant witnesses must survive admissible transport. The electromagnetic skin exposes only one invariant witness, $e$. Preserving independence would require loss of tensor load, hidden unwitnessed structure, or representational privilege. Each constitutes standing failure.

\section{Forced Tensor Relation}
The electromagnetic tensor load is uniquely specified by
\begin{equation}
\boxed{\;e = g_2 \sin\theta_W = g_1 \cos\theta_W,\qquad \tan\theta_W = \frac{g_1}{g_2}\;}
\end{equation}
For the electromagnetic sector, $g_2$, $g_1$, and $\theta_W$ are not independently admissible tensor degrees once $e$ is fixed.

\section{Corollaries}
\begin{itemize}
\item The electromagnetic sector admits a one-dimensional tensor load.
\item Electroweak and electromagnetic descriptions are tensor-identical for EM interactions after redundancy elimination.
\item Any framework asserting two independently tunable electromagnetic invariant couplings violates standing unless it introduces forbidden tensor load.
\end{itemize}

\section*{Conclusion}
Equivalence completion enforces collapse of apparent freedom in electroweak coupling bookkeeping for the electromagnetic sector. The result is eliminative, representation-invariant, and irreversible under admissible redescription.

\end{document}

